\section{Cross-repeat covariance: a point-estimate sensitivity}
\label{app:repeat-covariance}

This retrospective extension quantifies an assumed source of bias in the
primary 31-feature, full-frame error. It does not estimate the actual service
covariance, change the original uncertainty intervals, or add significance
tests. All six configurations and 15 pairwise point comparisons are retained.
The separate covariance-weighted \emph{feature metric} in
Appendix~\ref{app:magnitude-coverage} is a different diagnostic.

\subsection{Assumptions and quantities}
Stack the four artist-centered vectors for model $m$, scene $s$ and repeat $k$
as $d_{msk}=\theta_{ms}+e_{msk}$. Assume the same conditional mean in both
repeats, zero-mean errors and finite second moments. Define their average
normalized marginal energy and cross-repeat covariance trace as
\begin{equation}
 v_m=\frac{1}{SH}\sum_s\frac{\mathbb E\|e_{ms1}\|^2+
 \mathbb E\|e_{ms2}\|^2}{2},\qquad
 c_m=\frac{1}{SH}\sum_s\mathbb E\langle e_{ms1},e_{ms2}\rangle.
\end{equation}
Here $H$ is the original reference-contrast energy. Centering already accounts
for dependence among artists within a repeat. The exact expectation identities
are
\begin{equation}
 \mathbb E\widehat D_m=D_{*,m}+c_m,\qquad
 \widehat q_m=\frac{1}{2SH}\sum_s\|d_{ms1}-d_{ms2}\|^2,
 \qquad \mathbb E\widehat q_m=v_m-c_m,
\end{equation}
where $D_{*,m}$ is squared error of the stable scene-conditional mean contrasts.
Thus the observed difference statistic estimates marginal noise energy only
when $c_m=0$. These expectation identities do not require independent scenes
or equal marginal repeat variances; those assumptions matter for precision.
If repeat means differ, the difference statistic also contains their squared
separation and this stable-mean analysis does not apply.

For $v_m>0$, define the trace ratio $\rho_m=c_m/v_m\in[-1,1]$, called trace
correlation here. It is not an estimated Pearson correlation of observed
image vectors. Under an assumed common value $\rho_m=\rho<1$,
\begin{equation}
 \alpha(\rho)=\frac{\rho}{1-\rho},\qquad
 \widehat D_m(\rho)=\widehat D_m-\alpha(\rho)\widehat q_m.
 \label{eq:covariance-scenario}
\end{equation}
Equal trace correlation allows different covariance biases because models
have different noise scales. Equal normalized covariance $c_m=c$ would instead
shift every model equally and leave all pairwise differences unchanged. The
common-$\rho$ scenario does not cover arbitrary model-specific dependence.

\subsection{Fixed scenarios and point-order crossings}
The grid $\rho\in\{0,.10,.25,.50,.75\}$ was fixed before computing these
scenario outcomes. Negative adjusted estimates are retained; they are not
negative physical squared errors or evidence that the assumption has been
statistically rejected. No original interval is shifted or recomputed.
For each pair, let $\Delta_D=\widehat D_a-\widehat D_b$ and
$\Delta_q=\widehat q_a-\widehat q_b$. An interior positive crossing exists
when $\Delta_q\ne0$ and $\alpha_*:=\Delta_D/\Delta_q>0$, at
\begin{equation}
 \rho_* = \frac{\alpha_*}{1+\alpha_*}.
\end{equation}
The threshold describes these plug-in point curves; it is not a bound on actual
correlation, a confidence interval or a threshold for statistical superiority.
The complete tables below retain every pair and distinguish absent crossings
and ties. Unrounded values determine all calculations.

\subsection{What cannot be learned from two repeats}
Under this second-moment model, repeat differences alone impose no finite
upper bound on positive cross-repeat covariance. A random state in the centered
painter contrasts, shared by both repeats, cancels from their difference. Writing $q_*=\mathbb E\widehat q=v-c$,
for positive $q_*$, letting $\rho\uparrow1$ makes
$c=\rho q_* /(1-\rho)$ arbitrarily large. A zero observed
difference is also compatible with a variable shared state. Negative covariance
is possible and is outside this positive grid. Randomness over service sessions
and conditioning on one session's realized state also define different targets;
this collection does not distinguish them.

The scene-averaged shared naming fractions need separate common/baseline noise
and cross-scene covariance terms. Primary $q$ does not identify them, so no
fraction correction is inferred from this sensitivity. Likewise, the linear
aligned response remains unbiased under a mean-zero covariance model, but its
interval coverage is not validated here. Independently interleaved repeated
collections would provide information that these assumed curves cannot.
